\section{Piaci kontextus}

A WatchDB termékkoncepciójának értelmezéséhez érdemes röviden áttekinteni azokat a meglévő megoldásokat, amelyek részben hasonló problémára adnak választ. A piaci kontextus célja ebben az esetben nem egy teljes versenytárselemzés, hanem annak bemutatása, hogy a filmek követése, listázása és megosztása már létező felhasználói igény, ugyanakkor a különböző szolgáltatások eltérő megoldásokkal közelítik meg ezt a problémát.

Az IMDb egyik fontos releváns funkciója a watchlist. A felhasználók címeket adhatnak a saját watchlistjükhöz, rendezhetik azt, illetve értékelések és listák alapján is fedezhetnek fel új tartalmakat \cite{imdbWatchlistFaq}. Ez a megoldás jól mutatja, hogy a filmek későbbi megtekintésre való elmentése alapvető igény. A WatchDB szempontjából ugyanakkor az IMDb inkább információközpontú példaként jelenik meg, míg a projekt célja egy letisztultabb, személyes filmnapló és később közösségi élmény köré szerveződő alkalmazás.

A Trakt erősebb követési és integrációs szemléletet képvisel. A szolgáltatás filmek és sorozatok követésére, listák kezelésére, ajánlásokra, valamint különböző alkalmazásokkal és médiaközpontokkal való összekapcsolásra épít \cite{traktDiscoverTrackShare}. Ez a megközelítés funkcionálisan erős, ugyanakkor összetettebb termékélményt eredményezhet. A WatchDB PRD-je ezért inkább egy egyszerűbb, filmközpontú induló verziót határoz meg, ahol az elsődleges cél nem a sok integráció, hanem a gyors filmkeresés, naplózás és listaépítés.

A piacon jelen vannak kimondottan közösségi filmnapló platformok is, például a Letterboxd. A Letterboxd a filmek nyilvántartását, értékelését, naplózását, listázását és a barátok követését is támogatja, vagyis erősen kapcsolódik ahhoz a közösségi irányhoz, amely a WatchDB hosszabb távú víziójában is megjelenik \cite{letterboxdSocialFilmDiscovery}. Ez azt mutatja, hogy a filmnapló és a közösségi filmfelfedezés kombinációja valós felhasználói igényre épül. A WatchDB dolgozatbeli szerepe azonban nem az, hogy minden létező hasonló szolgáltatással teljes funkcionalitásban versenyezzen, hanem az, hogy egy ilyen termék tervezési és fejlesztési folyamatát egy szűkített MVP-n keresztül mutassa be.

A PRD-ben szereplő piaci összehasonlítás alapján a WatchDB lehetősége abban áll, hogy a fenti irányokból egy tudatosan szűkített változatot határoz meg. Az IMDb adatgazdagságából, a Trakt követési szemléletéből, a streaming szolgáltatások kényelmes listakezeléséből és a Letterboxd közösségi megközelítéséből is levonhatók tanulságok. Az MVP szempontjából azonban ezek közül azok a funkciók kapnak elsőbbséget, amelyek a legszélesebb célközönség számára azonnali értéket adnak: filmkeresés, watchlist, megnézett filmek naplózása és egyszerű értékelés. A közösségi és haladó funkciók így nem elvetett elemek, hanem a termék későbbi fejlődési irányai.
